\section{Localization, germs, and dense subcategories}
\label{sec:localization}

We will identify quotient morphisms with germs near the complement of
an open set, obtain the Witt localization sequence, and compare the
quotients arising from an open cover.  This open-cover quotient
comparison has kernel and cokernel annihilated by $2$; this integral
bound will control the denominators in the sphere calculation.

Throughout this section, $\Bcat$ denotes either $\T$ or $\Pcat$, and
$\Bzero$ denotes its triangulated subcategory before adjoining summands.
All spaces and, for $\Pcat$, all maps are as in
\Cref{def:test-categories}.  An open subset of a semialgebraic model need
not itself be semialgebraic.  The constructions below cut polyhedral
sources inside its inverse image, so this causes no additional assumption.
We use the compact support bounds, dualities, and fully faithful open
extension functors established in \Cref{prop:test-duality,prop:test-pushforward}.
We identify $\Bcat(U)$ with its image in $\Bcat(X)$ under extension by zero
when $U\subset X$ is open.

\subsection{Cutting off a finite construction}

The first step is to replace a finite construction by one supported
inside a prescribed open set while preserving its germ along a closed
subset.  An analogous cutoff for categories generated by simplex maps
appears in Ranicki--Weiss's local chain-complex construction
\cite[author version, Lemma~4.14 and Corollary~4.15]{RanickiWeiss2010}.
We prove the required statement here for the generated sheaf categories.

\begin{lemma}[Cutoff]
\label{lem:cutoff}
Let $K\subset O\subset X$, where $K$ is closed in $X$ and $O$ is open.
For every $A\in\Bzero(X)$ there is a morphism
\[
 s:A'\longrightarrow A
\]
such that $A'$ is extended from $\Bzero(O)$, its cohomology vanishes
outside a compact subset of $O$, and $s$ restricts to an isomorphism on
some open neighborhood of $K$.
\end{lemma}

\begin{proof}
First take $A=Rf_*\const{\R}{\Delta}$, with $\Delta$ a compact simplex.
The compact set $F=f^{-1}K$ lies in the relatively open set $f^{-1}O$.
Finite subdivision gives a relatively open polyhedral neighborhood $E$
of $F$ such that
\[
 C=\overline E\subset f^{-1}O,
 \qquad C\text{ and }C\setminus E\text{ are compact finite polyhedra}.
\]
For completeness, choose the subdivision so that the closed stars of
vertices of simplices meeting $F$ all lie in $f^{-1}O$; the positive
distance from $F$ to the complement and a sufficiently small mesh
ensure this.  Assign value zero to those vertices and value one to
all other vertices, and extend linearly to a function $h$.  It is zero
on $F$ and one outside $f^{-1}O$.  Take $E=\{h<1/2\}$ and
$C=\{h\leq1/2\}$, subdividing along the level set.  Because all vertex
values are zero or one, $C=\overline E$.
If $F$ is empty, take $E=C=\varnothing$.

For $e:E\hookrightarrow\Delta$, the natural map
$e_!\const{\R}{E}\to\const{\R}{\Delta}$ induces
\[
 s:A'=Rf_*e_!\const{\R}{E}\longrightarrow A.
\]
On $C$, the open--closed triangle is
\[
 e_{C!}\const{\R}{E}\longrightarrow\const{\R}{C}
 \longrightarrow i_*\const{\R}{C\setminus E}\longrightarrow
 e_{C!}\const{\R}{E}[1],
\]
where $e_C:E\hookrightarrow C$ and $i:C\setminus E\hookrightarrow C$.
Push this triangle along the restrictions of $f$.  Both compact sources
map into $O$, and therefore $A'$ is extended from $\Bzero(O)$, with
support bounded by the compact set $f(C)\subset O$.  In the
semialgebraic case these source subsets and restricted maps are still
semialgebraic.  Moreover,
\[
 N=X\setminus f(\Delta\setminus E)
\]
is open and contains $K$.  Its inverse image lies in $E$, so open
restriction of proper direct image identifies $s|_N$ with an
isomorphism.

The assertion is preserved by shifts and finite sums.  We check closure
under cones, including the compatibility of the morphisms.  Suppose
$a:A\to B$ is a morphism between objects for which the assertion holds.
First choose a cutoff $b:B'\to B$ over $O$, invertible on an open
$N\supset K$.  Cut off $A$ over the smaller open set $O\cap N$, obtaining
$c:A'\to A$.  Since $A'$ is supported in a compact subset of $N$, for
$j:N\hookrightarrow X$ we have $A'\simeq j_!j^*A'$.  The adjunction
\[
 \Hom_{\mathsf D(X)}(A',B')
 \simeq \Hom_{\mathsf D(N)}(A'|_N,B'|_N)
\]
therefore lifts
$(b|_N)^{-1}(a|_N)(c|_N)$ to a morphism $a':A'\to B'$.
The same adjunction with target $B$ proves the equality
$ba'=ac$ on $X$.  Complete this commutative square to a morphism of
distinguished triangles.  It gives
\[
 \Cone(a')\longrightarrow\Cone(a)
\]
which is invertible on the intersection of the neighborhoods where $b$
and $c$ are invertible.  Full faithfulness of open extension identifies
$a'$ with a morphism in $\Bzero(O)$, so its cone belongs to that category.
The union of the compact support bounds for $A'$ and $B'$ is a compact
support bound inside $O$ for this cone.  The finite triangulated
generation of $\Bzero(X)$ now proves the assertion.
\end{proof}

Cutoff also recognizes which summands come from an open set.

\begin{proposition}[Objects supported in an open set]
\label{prop:compact-support}
For an open inclusion $j:U\hookrightarrow X$, the essential image of
$j_!: \Bcat(U)\to\Bcat(X)$ consists exactly of those objects whose
cohomology vanishes outside some compact subset of $U$.  This image is
a thick, duality-stable subcategory.
\end{proposition}

\begin{proof}
An extended object has such a support bound by
\Cref{prop:test-pushforward}.  Conversely, let $S\in\Bcat(X)$ vanish
off a compact set $C\subset U$.  Write it as an actual summand
\[
 S\xrightarrow{i}A\xrightarrow{r}S,
 \qquad ri=\id_S,\qquad A\in\Bzero(X).
\]
By \Cref{lem:cutoff}, choose $s:A'\to A$ over $U$ which is invertible
on a neighborhood $N$ of $C$.  Since $S$ is extended by zero from $N$,
open-extension adjunction lifts $i$ to $i':S\to A'$ and gives
$si'=i$ globally.  Consequently $(rs)i'=\id_S$.
Restricting this splitting to $U$ makes $S|_U$ a summand of the object
$A'|_U\in\Bzero(U)$.  Thus $S|_U\in\Bcat(U)$, and its extension is $S$.

Finite unions of compact subsets of $U$ are compact in $U$.  The long
cohomology sequence of a triangle and the stalkwise splitting of a
summand show that this support condition is thick.  If $S$ vanishes
off $C$, open restriction of Verdier duality shows that $\D_X S$ also
vanishes off $C$.  Duality stability follows from
\Cref{prop:test-duality}.
\end{proof}

\subsection{The quotient as a category of germs}

The quotient by $\Bcat(U)$ forgets objects supported in compact subsets
of $U$.  We now show that its morphisms are determined by their
behavior on sufficiently small neighborhoods of the complement.

For $U\subset X$ open, put $K=X\setminus U$.  The neighborhoods of $K$
are ordered by shrinking; intersection makes this a filtered system.
For $A,B\in\Bcat(X)$, define the group of germ morphisms by
\begin{equation}
\label{loc:eq-germs}
 \operatorname*{colim}_{\substack{N\supset K\\N\text{ open in }X}}
 \Hom_{\mathsf D(N)}(A|_N,B|_N).
\end{equation}
Here $\mathsf D(N)=D^b(\operatorname{Sh}_{\R}(N))$ is the ambient
bounded derived category.
Neither restricted object is required to belong to
$\Bcat(N)$.  Composition is defined after restricting two representatives
to a common neighborhood.

\begin{theorem}[Germ description]
\label{thm:germs}
Restriction induces natural isomorphisms
\[
 \Hom_{\Bcat(X)/\Bcat(U)}(A,B)
 \simeq
 \operatorname*{colim}_{N\supset X\setminus U}
 \Hom_{\mathsf D(N)}(A|_N,B|_N)
\]
for every $A,B\in\Bcat(X)$.  These isomorphisms respect composition,
open restriction of duality, and restriction of distinguished triangles.
The quotient here is the Verdier quotient itself.
\end{theorem}

\begin{proof}
An object of $\Bcat(U)$ vanishes outside a compact $C\subset U$, and
therefore vanishes on the neighborhood $X\setminus C$ of $K$.
Consequently a morphism with cone in $\Bcat(U)$ becomes invertible in
germs.  Restriction thus factors through the Verdier quotient.

We prove fullness first for $A\in\Bzero(X)$.  Represent a germ by
$v:A|_N\to B|_N$ on $N\supset K$.  Apply \Cref{lem:cutoff} with
$O=N$ to obtain $s:A'\to A$, where $A'$ is compactly supported inside
$N$ and $s$ is invertible near $K$.  Adjunction extends
$v(s|_N)$ to a morphism $t:A'\to B$ on $X$.  The cone of $s$ has a
compact support bound $C$ and vanishes on some open $M\supset K$.
It is therefore supported in the compact subset $C\setminus M\subset U$;
\Cref{prop:compact-support} puts it in $\Bcat(U)$.  The roof
\[
 A\xleftarrow{s}A'\xrightarrow{t}B
\]
defines a quotient morphism with germ $v$.
For a general $A\in\Bcat(X)$, choose
$A\xrightarrow{i}A_0\xrightarrow{p}A$ with $pi=\id_A$ and
$A_0\in\Bzero(X)$.  Lift the germ $v(p|_N)$ by the case just proved,
then compose its lift with $i$.  This proves fullness for arbitrary
summand sources and arbitrary targets.

For faithfulness, first let $t:A\to B$ be an actual morphism whose germ
is zero.  If $A\in\Bzero(X)$, take a neighborhood $N\supset K$ on
which $t$ is zero and a cutoff $s:A'\to A$ supported inside $N$.
Open-extension adjunction gives $ts=0$ globally.  Since $s$ becomes
invertible in the quotient, $t$ becomes zero there.  For a summand source
as above, apply this assertion to $tp:A_0\to B$ and use $t=(tp)i$.
Finally, every quotient morphism is represented by a roof
$A\xleftarrow{s}C\xrightarrow{t}B$ with cone of $s$ in $\Bcat(U)$.
Its denominator is invertible near $K$.  If the roof has zero germ,
the numerator $t$ has zero germ and hence is zero in the quotient by
the actual-morphism case.  This proves faithfulness for all roofs.

The functor used throughout was restriction.  Its compatibility with
composition, triangles, and open Verdier duality proves the remaining
statements; the cutoff choices were used only to prove that this
canonical functor is full and faithful.
\end{proof}

\subsection{Witt localization}

The support characterization supplies the kernel category for
localization.  Balmer--Walter's theorem then gives an exact sequence
of integral Witt groups in every degree.

\begin{theorem}[Witt localization]
\label{thm:witt-localization}
For every open $U\subset X$ and every $r\in\Z$, there is a natural
long exact sequence
\begin{equation}
\label{loc:eq-witt-localization}
 \cdots\longrightarrow W_r(\Bcat(U))
 \longrightarrow W_r(\Bcat(X))
 \longrightarrow W_r(\Bcat(X)/\Bcat(U))
 \xrightarrow{\partial}W_{r-1}(\Bcat(U))
 \longrightarrow\cdots.
\end{equation}
The sequence is natural for $\Pcat\to\T$ and for the pushforwards
of \Cref{prop:test-pushforward}: a map $g:X\to X'$ with
$g(U)\subset U'$ induces a morphism from the sequence for $(X,U)$
to the sequence for $(X',U')$.
\end{theorem}

\begin{proof}
We apply \cite[Theorem~2.1, pp.~132--134]{BalmerWalter2002}.
Let $\mathcal J=\Bcat(U)$ and $\mathcal K=\Bcat(X)$.
By \Cref{prop:compact-support}, $\mathcal J$ is a full thick
subcategory stable under the exact duality.  It is therefore a
saturated full triangulated subcategory in the terminology of that
theorem.

These categories are essentially small: compact finite polyhedra and
their maps to a fixed target have a set of representatives; finite
triangulated constructions use sets of derived morphisms; and adding
summands still gives a set of isomorphism classes.  We may thus work
with small equivalent categories.  Morphism groups in $\mathcal K$
are real vector spaces, so multiplication by $2$ is bijective.
The enriched
octahedral axiom $(\mathrm{TR4}^{+})$ holds in derived categories of
exact categories and passes to full triangulated subcategories and
localizations \cite[\S2, p.~133]{BalmerWalter2002}.  Hence it holds
for $\mathcal K\subset D^b(\operatorname{Sh}_{\R}(X))$.

Let $S$ be the class of morphisms whose cones belong to $\mathcal J$.
The ordinary Verdier quotient
$S^{-1}\mathcal K=\mathcal K/\mathcal J$ has full kernel
$\mathcal J$, and the duality preserves $S$: dualizing a cone
triangle gives a shift of the corresponding dual cone.  All the
localization hypotheses are now satisfied.  Applying the theorem
to every shifted duality and rewriting its cohomological indices as
$W_r=W^{-r}$ gives
\eqref{loc:eq-witt-localization}.

For the naturality assertion, $\Pcat\to\T$ preserves the support
subcategories and the duality.  Similarly, if $g(U)\subset U'$, a
compact support bound $C\subset U$ maps to the compact support bound
$g(C)\subset U'$.  The exact duality-preserving pushforward therefore
preserves the kernel subcategories and induces a functor of their
localizations.

To check the boundary map, use its construction in
\cite[\S2, p.~133]{BalmerWalter2002}.  For the chosen shifted duality,
a symmetric form in the quotient is isometric to the image of a
symmetric morphism whose cone belongs to the kernel category.
Its boundary is the Witt class of a compatible symmetric cone.
An exact functor with the stated duality comparisons carries that
morphism and its symmetric-cone diagram to such a representative
in the target.  Both routes around the naturality square can
therefore use this same transported diagram.  Independence of the
representative gives equality of the boundary classes.  The remaining
maps commute by construction, proving the stated naturality.
\end{proof}

\subsection{Open covers and an integral density bound}

For an open cover, the two relevant quotients see the same germs.
Cutoff gives the following density statement, whose effect on Witt
groups we will then bound integrally.

\begin{proposition}[An open cover gives a dense inclusion]
\label{prop:open-cover}
If $X=U\cup V$ with $U,V$ open, extension by zero induces a fully
faithful exact functor preserving duality
\begin{equation}
\label{loc:eq-open-cover}
 \Bcat(V)/\Bcat(U\cap V)
 \longrightarrow\Bcat(X)/\Bcat(U).
\end{equation}
Its replete image is dense: every object of the target is an actual
summand of an object in the image.  These functors commute with
$\Pcat\to\T$.
\end{proposition}

\begin{proof}
Extension from $V$ sends $\Bcat(U\cap V)$ into $\Bcat(U)$, so it
induces the asserted exact functor with duality.  Put $K=X\setminus U$;
then $K\subset V$ and $K=V\setminus(U\cap V)$.  Open neighborhoods
of $K$ contained in $V$ are cofinal among its open neighborhoods in
$X$: intersect any given neighborhood with $V$.  Applying
\Cref{thm:germs} on $V$ and on $X$ therefore identifies the morphism
groups on the two sides of \eqref{loc:eq-open-cover}.  This proves
full faithfulness.

For $A_0\in\Bzero(X)$, apply \Cref{lem:cutoff} with $O=V$ and this
$K$.  Its cutoff $A'_0\to A_0$ has cone compactly supported inside
$U$, so it becomes an isomorphism in $\Bcat(X)/\Bcat(U)$.  The object
$A'_0$ comes from $\Bzero(V)$.  Every $A\in\Bcat(X)$ is a summand
of such an $A_0$, proving density.  Fullness and exactness show that
the replete image is triangulated: a morphism in the image lifts to
the source, and its cone is the image of a cone there.  It is also
duality stable.  All functors used are the same open extensions for
$\Pcat$ and $\T$, which proves the final assertion.
\end{proof}

The exponent-two kernel and cokernel bound for dense inclusions follows
from the cofinality theorem for triangular Witt groups of
Hornbostel--Schlichting
\cite[author preprint, Appendix~A, Theorem~A.2]{HornbostelSchlichting2004}.
We give a direct proof with an explicit homomorphism in the reverse
direction.  It doubles forms into the smaller category and checks that
neutral relations descend there.

\begin{theorem}[Dense doubling]
\label{thm:dense-doubling}
Let $\mathcal A\subset\mathcal C$ be a full replete triangulated
subcategory stable under duality, and suppose it is dense: each object
of $\mathcal C$ is a summand of an object of $\mathcal A$.
Assume these categories admit triangular Witt groups.  In every degree,
the inclusion homomorphism
\[
 \iota:W_r(\mathcal A)\longrightarrow W_r(\mathcal C)
\]
has kernel and cokernel annihilated by $2$.  More precisely, there is
a homomorphism
$d:W_r(\mathcal C)\to W_r(\mathcal A)$ satisfying
\[
 \iota d=2\id,
 \qquad d\iota=2\id.
\]
No closure of $\mathcal A$ under arbitrary summands is required.
\end{theorem}

\begin{proof}
Fix the degree and denote its duality by $\#$.  We first record a
consequence of density:
\begin{equation}
\label{loc:eq-adjacent-shifts}
 y\oplus y[1]\in\mathcal A
 \qquad(y\in\mathcal C).
\end{equation}
Indeed, write $c\simeq y\oplus z$ with $c\in\mathcal A$.
If density is expressed by a split inclusion, its cone supplies such
a complement by the split-triangle identity.  The endomorphism
$0_y\oplus\id_z$ of $c$ belongs to $\mathcal A$ by fullness.
Its cone is $y\oplus y[1]$, which proves
\eqref{loc:eq-adjacent-shifts}.  Applying this to $y[-1]$ also gives
$y\oplus y[-1]\in\mathcal A$.

For a nonsingular symmetric form $q=(x,\alpha)$ in $\mathcal C$, set
\begin{equation}
\label{loc:eq-doubling-form}
 D(q)=q\perp q\perp H(x[1]),
\end{equation}
where $H$ denotes the hyperbolic form for $\#$.  Its underlying object
is
\[
 x\oplus x\oplus x[1]\oplus\#(x[1])
 \simeq (x\oplus x[1])\oplus(x\oplus x[-1]),
\]
using $\alpha:x\simeq\#x$ and
$\#(x[1])\simeq(\#x)[-1]$.
By \eqref{loc:eq-adjacent-shifts} it belongs to $\mathcal A$.
Repleteness and fullness make $D(q)$ a form in $\mathcal A$.
This assignment respects isometries and orthogonal sums: the direct
sum of isometries acts on the two copies of $q$, and the hyperbolic
construction respects isomorphisms and direct sums.  It remains to
check neutral forms.

Suppose $q$ has a triangular lagrangian $a:l\to x$.  Its lagrangian
triangle has the form
\[
 l\xrightarrow{a}x\longrightarrow\#l\longrightarrow l[1],
\]
with the duality compatibility that defines a neutral form.
The object $z=l\oplus l[1]$ belongs to $\mathcal A$, so $H(z)$ is
neutral already in $\mathcal A$.  In $\mathcal C$, the form
$D(q)\perp H(z)$ has lagrangian
\begin{equation}
\label{loc:eq-doubled-lagrangian}
 L=l\oplus l\oplus x[1]\oplus\#l\oplus l[1].
\end{equation}
Here the first two factors are the lagrangians of the copies of $q$;
$x[1]$ is the first lagrangian of $H(x[1])$; and, in
$H(z)\simeq H(l)\perp H(l[1])$, choose the second lagrangian $\#l$
of $H(l)$ and the first lagrangian $l[1]$ of $H(l[1])$.
Orthogonal sum of these lagrangian triangles gives the indicated
duality-compatible lagrangian triangle.

We verify $L\in\mathcal A$ using only fullness and triangles.
Both objects of the morphism
\[
 a\oplus0:l\oplus l[1]\longrightarrow x\oplus x[1]
\]
belong to $\mathcal A$ by \eqref{loc:eq-adjacent-shifts}.  Its cone,
and hence the following object, belongs to $\mathcal A$:
\[
 C=\#l\oplus x[1]\oplus l[2].
\]
Add $l\oplus l[1]\in\mathcal A$ and rearrange to obtain
\[
 C\oplus l\oplus l[1]
 \simeq
 (\#l\oplus x[1]\oplus l)\oplus(l[1]\oplus l[2]).
\]
The second displayed summand $l[1]\oplus l[2]$ is itself in
$\mathcal A$.  Its split inclusion into the displayed total object
is therefore a morphism in $\mathcal A$; taking its cone puts
$\#l\oplus x[1]\oplus l$ in $\mathcal A$.  Finally,
\[
 L\simeq(l\oplus l[1])\oplus(\#l\oplus x[1]\oplus l)
\]
also belongs to $\mathcal A$.  This removal used a split triangle
between objects already in $\mathcal A$.

The underlying object of $D(q)\perp H(z)$, its lagrangian $L$, and
$\#L$ now belong to $\mathcal A$.  Fullness supplies all arrows of
the lagrangian diagram in $\mathcal A$; its inherited triangulation
and duality make the same diagram a lagrangian diagram there.
Thus $D(q)\perp H(z)$ is neutral in $\mathcal A$, and since $H(z)$
is neutral there, $[D(q)]=0$ in $W_r(\mathcal A)$.
The assignment \eqref{loc:eq-doubling-form} consequently descends to
a homomorphism $d:W_r(\mathcal C)\to W_r(\mathcal A)$.

In $\mathcal C$, $H(x[1])$ is neutral, so
$\iota d[q]=2[q]$.  If $q$ starts in $\mathcal A$, then $x[1]$ is
in $\mathcal A$ and the same hyperbolic form is neutral there, giving
$d\iota[q]=2[q]$.  The first identity annihilates the cokernel of
$\iota$ by $2$, and the second annihilates its kernel by $2$.
\end{proof}

\begin{corollary}
\label{loc:cor-open-cover-bound}
The map on Witt groups induced by \eqref{loc:eq-open-cover} has kernel
and cokernel annihilated by $2$, in every degree and for either test
category.
\end{corollary}

\begin{proof}
Apply \Cref{thm:dense-doubling} to the replete image in
\Cref{prop:open-cover}.
\end{proof}
